% !TeX root = ../fundamentals.tex
% Section 2.3 -- Multi-task learning
% Draft 1 (2026-07-21). Obeys WRITING_LAW.md + GLOSSARY.md. Citation/number ledger at the section end.

\section{Multi-task learning}
\label{sec:fund:mtl}

Multi-task learning (MTL) trains one model on several related tasks at once, in the
expectation that a representation shared among them generalizes better than one
learned for a single task in isolation \cite{caruana1997multitask}. The appeal for
the present work is direct: predicting the next category and the next region from the
same check-in history may let each task inform the other. Whether it does is the
question the dissertation opens with, and the answer turns out to depend on choices
this section lays out.

Deep MTL is organized around how much the tasks share. The two poles are hard
parameter sharing, where the tasks use a common trunk and split only at the output
heads, and soft parameter sharing, where each task keeps its own parameters and the
tasks are coupled by a weaker constraint \cite{ruder2017mtloverview}. Between the
poles lie architectures that learn what to share. Cross-stitch units place a learned
linear combination of two task networks' activations at each layer, letting the model
decide how much to mix them \cite{misra2016cross}. The multi-gate mixture-of-experts
shares a pool of expert subnetworks but gives each task its own gate over them, so
that routing, not a fixed trunk, determines sharing \cite{ma2018mmoe}. Progressive
layered extraction sharpens this by separating shared experts from task-specific ones
in layered gates \cite{tang2020ple}, and DSelect-k makes the expert selection itself
differentiable and sparse \cite{hazimeh2021dselectk}. The principle these architectures
share, keeping shared and task-specific components side by side rather than forcing one
common trunk, is the one the joint model of Chapter~5 adopts, though it realizes it with
cross-attention rather than expert gating.

Sharing is not free. When tasks pull the shared parameters in different directions,
joint training can leave a task worse off than its single-task model, a failure
known as negative transfer. Which tasks help each other and which interfere cannot be
assumed away, and joint training can hurt as easily as it helps depending on the
pairing \cite{standley2020tasks}. The effect is not incidental: because the tasks can
conflict, minimizing a weighted sum of their losses is only valid when they do not,
and multi-task learning is more honestly cast as a multi-objective problem with no
single optimum \cite{sener2018mgda}.

A family of methods tries to manage the conflict at the level of the gradients or the
losses. Uncertainty weighting sets each task's loss weight from its homoscedastic
uncertainty, learned jointly with the model \cite{kendall2018uncertainty}. GradNorm
rescales per-task gradient magnitudes so that tasks train at comparable rates
\cite{chen2018gradnorm}, and dynamic weight averaging sets the weights from the recent
rate of change of each loss \cite{liu2019dwa}. Other methods act on the gradient
directions: PCGrad projects a task's gradient off any conflicting task's gradient
\cite{yu2020pcgrad}, CAGrad seeks an update close to the average gradient while
bounding conflict, with convergence guarantees \cite{liu2021cagrad}, Nash-MTL treats
the combination of task gradients as a bargaining game and takes the Nash solution as
the joint direction \cite{navon2022nashmtl}, and Aligned-MTL aligns the principal
components of the gradient system and uses its condition number as a stability
criterion \cite{senushkin2023aligned}. FAMO reduces the cost of balancing by tracking
loss decreases in constant time and space rather than reading every task gradient at
each step \cite{liu2023famo}. A recurring result tempers the family. Random loss
weighting is a competitive baseline \cite{lin2022rlw}; a controlled study finds that
current MTL optimizers often do not outperform a well-tuned fixed-weight baseline
\cite{xin2022domtl}; and a direct defense of plain loss summation shows that a unitary
scalarization with standard regularization matches or improves upon the specialized
optimizers \cite{kurin2022scalarization}. A recent survey and a dense-prediction survey
both consolidate the architectures and the balancing methods and set out when sharing
helps \cite{vandenhende2022mtl,yu2024survey}. The dissertation takes the cautious
position these results support: a fixed-weight baseline is a serious competitor, and a
balancer earns its place only by outperforming it.

In mobility, MTL has been used almost entirely in the service of next place. MCARNN
predicts activity and next place together with a context-aware recurrent unit, pairing
the two in parallel \cite{Liao2018}. A cascade line runs the other way: CSLSL imposes
a causal chain, time then activity then location, so that the activity or category is
an intermediate step toward the place \cite{huang2024cslsl}. Across these systems the
region and the category appear as auxiliary signals or as stages on the path to a
place; no multi-task model among them predicts the next region as a co-equal end
target alongside the next category. Single-task work on region prediction exists,
as Section~\ref{sec:fund:tasks} notes, but the joint setting is open. That absence
is the opening the dissertation addresses. Its own starting model, MTLnet, applies
hard parameter sharing on a place-level embedding and does not outperform the dedicated
single-task models, a result that holds for that configuration and motivates the
representation argument the next chapters develop \cite{silva2025mtlnet}.

% ---------------------------------------------------------------------------
% CITATION / NUMBER LEDGER -- 2.3 (verbs bound to tests; no result upgraded)
%  caruana1997multitask   MTL = inductive transfer via shared representation (defines MTL)
%  ruder2017mtloverview   hard vs soft parameter sharing (frames design space) [new bib, @misc arXiv:1706.05098]
%  misra2016cross         cross-stitch soft sharing [ERRATA DOI -> 10.1109/CVPR.2016.433]
%  ma2018mmoe             MMoE: per-task gates over shared experts (routing)
%  tang2020ple            PLE/CGC: shared vs task-specific experts (structured sharing) [new bib]
%  hazimeh2021dselectk    DSelect-k: differentiable sparse expert selection [new bib]
%  PLE lineage FIX (domain F6): joint model adopts the structured-sharing PRINCIPLE (shared+task-specific),
%                        realized with cross-attention, NOT a PLE/MoE descendant. Reworded, no "descends from".
%  standley2020tasks      negative transfer / which tasks to learn together; joint can hurt (VERIFIED arXiv abstract; existing CBIC key).
%                        FIX (both gates F2): negative transfer DEFINED here, folded onto standley -- Zhang2020 REMOVED
%                        (Zhang2020 = iMTL next-POI recommender, NOT negative-transfer literature; reserved for mobility use if needed).
%  sener2018mgda          MTL as multi-objective optimization; weighted sum valid only if no conflict (VERIFIED arXiv abstract; existing CBIC key)
%  kendall2018uncertainty homoscedastic-uncertainty loss weighting [new bib; DOI 10.1109/CVPR.2018.00781 verified + arXiv:1705.07115]
%  chen2018gradnorm       GradNorm: gradient-magnitude balancing
%  liu2019dwa             DWA: loss-rate schedule weights
%  yu2020pcgrad           PCGrad: gradient projection off conflict
%  liu2021cagrad          CAGrad: conflict-averse update w/ convergence [new bib]
%  navon2022nashmtl       Nash-MTL: bargaining-game gradient combination [ERRATA: consolidate nash double-key/slash]
%  senushkin2023aligned   Aligned-MTL: principal-component alignment, condition number
%  liu2023famo            FAMO: O(1) loss-decrease balancing
%  lin2022rlw             RLW: random weighting is competitive [new bib]
%  xin2022domtl           controlled study: MTL optimizers often do not outperform tuned fixed weights
%  kurin2022scalarization Kurin et al. NeurIPS 2022 "In Defense of the Unitary Scalarization": plain loss-sum + std
%                        regularization matches/improves specialized optimizers (VERIFIED firsthand arXiv 2201.04122 abstract
%                        + NeurIPS 2022 venue). Added per domain F4 (canonical scalarization-skeptic anchor). [new bib]
%  vandenhende2022mtl     dense-prediction MTL survey (anchor)
%  yu2024survey           recent comprehensive MTL survey (existing CBIC key)
%  Liao2018               MCARNN: activity + next place (parallel MTL, mobility)
%  huang2024cslsl         CSLSL: causal cascade time->activity->location (category intermediate)
%  silva2025mtlnet        MTLnet null result: hard sharing on place embedding ~ single-task (TIME-INDEXED, "for that configuration")
%  NUMBERS: none quoted. MTLnet null stated as configuration-specific ("does not outperform ... for that configuration"), no metric upgraded.
%  HONESTY: banned "beat" REMOVED (domain/fact F7): all four instances -> "outperform"/reworded. Region-as-end-target claim
%           SCOPED to the multi-task co-equal setting (domain/fact F3): single-task region prediction (zhu2022drrgnn, 2.1) exists.
%           Null result written plainly and time-indexed; no result verb upgraded.
% ---------------------------------------------------------------------------
